% =====================================================================
%  Network Intrusion Detection on CICIDS2017
%  Springer LNCS / CCIS format  (target venues: ICTA 2026, FDSE 2026)
%  Compiled FREE on GitHub Actions via xu-cheng/latex-action@v3
%
%  GOLDEN RULE: every number in this paper must come from the Colab
%  experiments (experiments/ids_cicids2017.ipynb). Placeholders are
%  marked TODO and must be replaced with real measured values.
% =====================================================================
% SOICT requires the submitted PDF to carry no page numbers, so the
% running heads are dropped for submission.
\documentclass{llncs}
% SOICT 2026 is single-blind: \anonfalse keeps author information visible.
% Set \anontrue only for a double-blind venue.
\newif\ifanon\anonfalse

\usepackage{graphicx}
\usepackage{booktabs}
\usepackage{amsmath}
\usepackage{multirow}
\usepackage{array}
\usepackage{url}
\usepackage{tikz}
\usetikzlibrary{positioning,fit,arrows.meta}
% NOTE: hyperref intentionally NOT loaded — it clashes with llncs
% (\@xdblarg error). Springer adds hyperlinking for the final version.

% Tables are generated by experiments/make_tables.py directly from
% experiments/results_v3/*.json, so no measured number is typed by hand.
% A stage that has not been run yet simply contributes no table.
\newcommand{\inputtable}[1]{%
  \IfFileExists{tables/#1.tex}{\input{tables/#1}}{}}

% Figures are generated by experiments/make_figures.py from the same
% JSON as the tables. A figure whose stage has not run is skipped rather
% than breaking the build.
\newcommand{\resfig}[3]{%
  \IfFileExists{figures/#1.pdf}{%
    \begin{figure}[tb]\centering
    \includegraphics[width=0.88\textwidth]{figures/#1}
    \caption{#2}\label{#3}
    \end{figure}}{}}

% The paper carries many small floats; the default separation costs more
% than a page across the document.
\setlength{\textfloatsep}{9pt plus 2pt minus 2pt}
\setlength{\intextsep}{9pt plus 2pt minus 2pt}
\setlength{\floatsep}{7pt plus 2pt minus 2pt}
\renewcommand{\topfraction}{0.9}
\renewcommand{\bottomfraction}{0.7}
\renewcommand{\textfraction}{0.08}
\renewcommand{\floatpagefraction}{0.8}

\pagestyle{empty}
\begin{document}
\thispagestyle{empty}

\title{Network Intrusion Detection on CICIDS2017: A Comparative Study of
Machine Learning and Deep Learning Models with a Cloud-Native
Deployment Design}

\ifanon
\author{Anonymous Author(s)}
\institute{Affiliation withheld for double-blind review}
\else
\author{Huy-Thanh Pham\inst{1} \and Tuan-Anh Truong\inst{1} \and
Vinh-Khuong Truong\inst{2}}
\authorrunning{H.-T. Pham et al.}
\institute{Faculty of Computer Science and Engineering,\\
Ho Chi Minh City University of Technology (HCMUT), VNU-HCM, Vietnam \\
\email{huythanh@clouds.io.vn, anhtt@hcmut.edu.vn}
\and Ton Duc Thang University, Ho Chi Minh City, Vietnam \\
\email{trvinhkhuong@gmail.com}}
\fi

\maketitle

\begin{abstract}
Machine-learned intrusion detectors report near-perfect scores on public
benchmarks, and practitioners do not see that performance in the field. We
show that a large part of the discrepancy is produced by the evaluation
itself, and we measure how large. Holding the models, the hyperparameters and
the preprocessing fixed on CICIDS2017, and re-fitting feature selection,
oversampling and scaling inside every training split, we compare a
conventional stratified random split against a day-disjoint split in which no
attack family is shared between training and test. Attack recall for XGBoost
falls from 0.9935 to 0.3164 and macro-F1 from 0.9879 to 0.6887; Random Forest
behaves the same way. The ranking of models does not survive either: Logistic
Regression, the weakest detector under the random split, becomes the
strongest under the day-disjoint one, so a benchmark that flatters a model can
also select the wrong one. Repeating the conventional protocol over five
seeds further shows that the usual claim of a best model is unsupported --
XGBoost leads Random Forest by 0.00044 macro-F1, which neither a paired
$t$-test nor a Wilcoxon signed-rank test separates from zero. We then
translate the measured error rates into precision at operational attack
prevalence, where benchmark precision does not transfer, and we measure the model served as a
containerised endpoint under controlled load: it sustains 309~requests/s
before latency collapses, where the amortised per-sample model cost -- the
figure the literature, and an earlier version of this work, mistakes for a
service throughput -- would predict roughly $9.5\times10^{5}$~flows/s, an
overstatement of three orders of magnitude. The experiment
driver, the serving implementation, the deployment manifests and the raw
result files from which every table and figure here is generated are released
with the paper.
\keywords{Intrusion detection \and Evaluation methodology \and Data leakage
\and CICIDS2017 \and Model serving \and Network security.}
\end{abstract}

% =====================================================================
\section{Introduction}
\label{sec:intro}
Learning-based intrusion detection reports extraordinary accuracy on public
benchmarks -- macro-F1 above 0.98 is routine -- and practitioners do not
observe it in the field. The discrepancy is usually attributed to concept
drift or to the age of the corpora. This paper examines a more immediate
contributor: much of the reported accuracy is an artefact of how the
benchmark is split and of what the reported metric measures.

Two mechanisms are at work. First, a random split rewards recognition rather
than detection: flows from a single attack episode are numerous and
near-identical, and a random split scatters them across the train/test
boundary, so a model scores well by recognising near-duplicates it has
already memorised. That random splits flatter intrusion detectors has been
observed before~\cite{moczkodan2026transformers,xu2025robust} and we do not
claim the observation; we contribute a protocol that removes the confound
and a measurement of what it costs. Second, precision depends on the
prevalence of attacks, and benchmark corpora carry orders of magnitude more
attack traffic than a production link, so a precision measured at a 17.7\,\%
attack rate says little about the alert queue an analyst would face at one
in a thousand.

A third gap is one of engineering. Papers that propose a deployment
routinely support it with an \emph{amortised per-sample model cost} -- the
time to score a test batch divided by its size, on all cores -- which omits
feature extraction, parsing, serialisation, transport and concurrency.
It is a lower bound on the cost of the model, not a property of a service.
An earlier version of this work converted such a figure into a throughput,
which is what prompted the present study.

\paragraph{Contributions.}
We hold the models, hyperparameters and preprocessing fixed and vary only
the evaluation: (i) a pipeline in which scaling, feature selection and
oversampling are re-fitted \emph{inside} every training split
(Section~\ref{sec:leakage}); (ii) four splitting protocols on identical data
-- repeated holdout, day-disjoint (which on CICIDS2017 is also
family-disjoint), leave-one-family-out, and genuine cross-dataset transfer
(Section~\ref{sec:method}); (iii) paired statistical tests over repeated
seeds, which show the commonly reported ranking of the two leading models is
unsupported (Section~\ref{sec:paired}); (iv) a translation of measured error
rates into precision and alert burden at operational prevalence
(Section~\ref{sec:operational}); and (v) an end-to-end measurement of the
deployed service, reported separately from model cost
(Section~\ref{sec:serving}). The experiment driver, serving implementation,
container and Kubernetes definitions, load profile and raw result files --
from which every table and figure here is generated -- are released with the
paper.

\section{Related Work}
\label{sec:related}
Machine learning has been applied to intrusion detection for two decades; the
survey of Buczak and Guven~\cite{buczak2016survey} covers the method families
and their cost trade-offs. Our concern is narrower and is with evaluation
rather than models.

\paragraph{Benchmarks and their defects.}
NSL-KDD~\cite{tavallaee2009nslkdd} corrected redundant records and
unrealistic class balance in KDD'99; UNSW-NB15~\cite{moustafa2015unswnb15}
introduced a more contemporary feature set; CICIDS2017~\cite{sharafaldin2018cicids}
added five days of labelled traffic across several attack families.
CICIDS2017 is not a neutral instrument, however: Engelen et
al.~\cite{engelen2021troubleshooting} document labelling and capture defects
-- misordered and duplicated packets, executed but unlabelled attacks -- and
Lanvin et al.~\cite{lanvin2023errors} show that correcting them changes
measured performance materially. We therefore treat absolute scores as
properties of the benchmark and report differences between protocols on
identical data, which are far less sensitive to label noise.

\paragraph{Evaluation protocol.}
That random splits flatter intrusion detectors is not new and we do not claim
it. Moczkodan and Ragab~\cite{moczkodan2026transformers} compare random,
temporal and five-tuple-grouped splits on CICIDS2017 and find the random
convention substantially overestimates performance; Xu and
Liu~\cite{xu2025robust} report severe recall loss on novel attacks. Our
contribution extends the observation along three axes those studies leave
open: a \emph{day-disjoint} split, which on this corpus is simultaneously
family-disjoint; a translation of the resulting error rates into precision at
operational prevalence; and a measurement of what the chosen detector costs
to serve.

\paragraph{From model cost to service cost.}
That last axis is treated most loosely in the literature. It is common to
report an amortised per-sample model cost -- figures of order
$10^{-2}$~ms are typical -- as evidence of deployability. Such a figure
excludes feature extraction, parsing, serialisation, transport, queueing and
concurrency, and is not a property of a running service. We report it too,
but separately from, and never converted into, the end-to-end latency
distribution of a deployed endpoint. Table~\ref{tab:related} summarises the
positioning; recent applied studies on this
benchmark~\cite{talukder2024ml,almuhanna2025dl} confirm its continued use and
the strength of tabular models on it.

\inputtable{related_work}

\section{Data and Preprocessing}
\label{sec:data}
CICIDS2017~\cite{sharafaldin2018cicids} contains five working days of
labelled flows -- benign traffic plus denial-of-service, distributed
denial-of-service, brute-force, web-attack, botnet, port-scan and
infiltration families -- each described by more than eighty flow statistics,
distributed as eight CSV files, one per capture period. We draw a fixed
20\,\% stratified sample from each file. After normalising column names,
replacing infinities with missing markers, dropping incomplete records and
removing exact duplicates, the working set holds 535{,}069 flows with 78
numeric features and an attack ratio of 0.1769.

Two pieces of information a binary formulation would discard are retained,
because the protocols below depend on them: each flow keeps its
\emph{capture day}, recovered from its source file, which is what makes the
day-disjoint split possible; and its \emph{original multi-class label}, which
is what makes per-family analysis possible. The binary target maps every
non-benign label to the positive class.

Standardisation, top-$k$ selection by Random Forest
importance~\cite{breiman2001random} and oversampling with
SMOTE~\cite{chawla2002smote} are treated as part of the model rather than of
the data: each is fitted on the training portion of whichever split is in
force, never on the corpus as a whole (Section~\ref{sec:leakage}).

\section{Evaluation Protocol}
\label{sec:method}
The models in this study are deliberately ordinary. What varies, and what the
paper is about, is the protocol under which they are evaluated. This section
fixes that protocol precisely enough to be re-run.

\subsection{Models and Hyperparameters}
We compare four detectors spanning a range of capacity and cost: Logistic
Regression as a linear baseline, Random Forest~\cite{breiman2001random},
XGBoost~\cite{chen2016xgboost}, and a two-hidden-layer multilayer perceptron.

Table~\ref{tab:hyper} lists every hyperparameter. The values are fixed a
priori and identical across every split, seed and dataset; we perform no
tuning. The absolute scores are therefore not the best attainable, but no
model is advantaged by a larger search budget than another, and no
hyperparameter can have been selected -- however indirectly -- on test data.
Since the paper reports differences between protocols on a fixed pipeline
rather than a state-of-the-art claim, this removes a confound rather than
weakening the result. It is recorded as a limitation nonetheless.

\inputtable{hyperparameters}

\subsection{Leakage Control Inside the Pipeline}
\label{sec:leakage}
Preprocessing fitted before the split leaks test information into training
even when the classifier never sees a test row. Feature selection is the
worst offender: choosing the top-$k$ by importance computed over the whole
corpus lets the test set shape the input space of every later model.

We therefore treat standardisation, selection and oversampling as part of the
model. For every split in every protocol, the scaler is fitted on the
training portion only, the Random-Forest ranking selecting the top $k=20$
features is computed on the scaled training portion only, and
SMOTE~\cite{chawla2002smote} is applied to the training portion only, after
selection. The test portion is transformed by the fitted scaler and reduced
to the chosen features. Because the ranking is recomputed per split, the
selected set is itself a random variable, whose stability we report in
Section~\ref{sec:featstab} rather than presenting one list as canonical.

\subsection{Four Splitting Protocols}
The same pipeline is evaluated under four separations of training from test,
in increasing order of difficulty. Reporting them side by side is the point:
each answers a different question, and only the last two resemble what a
deployed sensor faces.

\textbf{(P1) Repeated stratified holdout} -- a 70/30 stratified split over
five seeds. This is the conventional protocol and the reference point.
Repetition is not cosmetic: one split cannot distinguish a real difference
between models from split-to-split noise.

\textbf{(P2) Day-disjoint} -- training on Monday--Wednesday, testing on
Thursday--Friday, so no capture day appears on both sides and the
near-duplicate flows a random split scatters across the boundary are removed.
On CICIDS2017 each attack family was executed on one specific day, so this
split is necessarily also \emph{family-disjoint}: every attack in the test set
is of a type absent from training, and P2 measures detection of unseen types
rather than recognition of familiar ones.

\textbf{(P3) Leave-one-family-out} -- P2 confounds the day boundary with the
family boundary; P3 separates them by removing one family from training
entirely and asking the model to flag it, family by family.

\textbf{(P4) Cross-dataset transfer} -- the strictest test changes corpus.
CICIDS2017 and UNSW-NB15 share no feature schema, so transfer requires an
explicit alignment: we use the seven flow-level quantities both provide with
compatible semantics, rescaling flow duration from microseconds to seconds. A
model is fitted on one corpus and evaluated on the other with no
target-domain data. We distinguish this from the weaker
\emph{independent-dataset replication}, in which the protocol is re-run within
UNSW-NB15 to check whether the model ordering reproduces. The two are
frequently conflated, including in an earlier version of this work; only P4
speaks to cross-domain behaviour.

\subsection{Statistical Treatment and Metrics}
\label{sec:stats}
Every protocol that admits repetition is repeated over multiple seeds, with
all models evaluated on identical splits, so comparisons are paired. We
report means with the half-width of the 95\,\% confidence interval, and for
the leading pair a paired $t$-test and a Wilcoxon signed-rank test over the
shared seeds, so that a claim of superiority rests on a test rather than on
a difference in the fourth decimal place.

Alongside accuracy, precision, recall and macro-F1 we report the false
positive and false negative rates, the Matthews correlation coefficient,
balanced accuracy, and the area under both the ROC and the
precision--recall curve; PR-AUC is the more informative of the two under
heavy imbalance, since ROC-AUC is insensitive to the absolute number of
false positives when negatives dominate.

\paragraph{Operational prevalence.}
Precision depends on prevalence while TPR and FPR do not, so a precision
measured on a benchmark does not transfer to deployment. For prevalence
$\pi$, the precision implied by a measured $(\mathrm{TPR}, \mathrm{FPR})$ is
\begin{equation}
\label{eq:ppv}
\mathrm{PPV}(\pi)=\frac{\pi\cdot\mathrm{TPR}}
{\pi\cdot\mathrm{TPR}+(1-\pi)\cdot\mathrm{FPR}},
\end{equation}
which we evaluate at $\pi \in \{10^{-1},10^{-2},10^{-3},10^{-4}\}$ together
with the number of false alerts per million flows -- the quantity that
decides whether a detector is usable at all.

\subsection{Cost Measurement}
Training time and inference time are measured on the environment of
Table~\ref{tab:env}. Inference is timed after a discarded warm-up pass, as
the mean of five repeats, on the full test batch with all cores available.
The resulting per-sample figure is an amortised lower bound on model cost.
We do not convert it into a service throughput. What a service actually
sustains is measured separately, on the running endpoint, in
Section~\ref{sec:serving}.

% =====================================================================
\section{Results}
\label{sec:results}
We report the four protocols of Section~\ref{sec:method} in order, then the
feature-selection analysis and the operational reading of the error rates.
Every number below is generated directly from the stored experiment output;
none is transcribed by hand.

% [trimmed for page limit]
%\inputtable{dataset}
% [trimmed for page limit] the two day groups are named in the text
%\inputtable{days}
\subsection{P1 and P2: Random versus Day-Disjoint}
\label{sec:p1}
Table~\ref{tab:protocols} reports both protocols on identical data. Under the
random split the two tree ensembles lead clearly, and the interval
half-widths deserve as much attention as the means: they are tight for Random
Forest and XGBoost and wide for Logistic Regression and the MLP
(Figure~\ref{fig:seedspread}), so a single-split number for either of the
latter is close to uninformative.

\inputtable{protocols}
% [trimmed for page limit] \resfig{fig_seed_spread}

\subsection{Is the Leading Model Actually Leading?}
\label{sec:paired}
The usual conclusion from a table like this is that XGBoost wins.
Table~\ref{tab:paired} tests the claim on the paired seeds instead of
asserting it. XGBoost does lead Random Forest, but by 0.00044 macro-F1, and
neither the paired $t$-test ($p=0.13$) nor the Wilcoxon signed-rank test
($p=0.19$) separates that from zero. Both ensembles are separated from
Logistic Regression and the MLP by margins that are large and clearly
significant.

We therefore do not name a best model. What the data supports is two tiers --
tree ensembles above the linear and shallow neural baselines -- with the
choice inside the upper tier resting on other grounds, which
Section~\ref{sec:cost} shows are not negligible.

\inputtable{paired_tests}

\subsection{P2: Day-Disjoint, and the Optimism of a Random Split}
\label{sec:p2}
Training uses Monday--Wednesday and testing Thursday--Friday. The two sides
share no attack family: training holds the denial-of-service variants, both
brute-force families and Heartbleed, while testing holds botnet, DDoS,
infiltration, port scanning and the three web-attack variants. Every attack
seen at test time is of a type absent from training.

The effect is not a third-decimal correction. Attack recall for XGBoost falls
from 0.9935 to 0.3164 and for Random Forest from 0.9937 to 0.3109; macro-F1
falls by roughly 0.30 for both. Under the conventional protocol these models
find essentially every attack in the test set; under this one they miss two
in three.

Figure~\ref{fig:confusion} shows where the errors go, and the direction
matters. The false-positive rate does not deteriorate -- it improves, from
0.0073 to 0.0030 -- while false negatives rise from 185 flows to 36{,}506.
The detector does not become noisy under the harder protocol; it becomes
quiet, confident and wrong, which is the failure a monitoring dashboard is
least likely to reveal.

Nor does the ranking survive. Logistic Regression is last under the random
split (macro-F1 0.8515 against 0.9879) and first under the day-disjoint one
(0.7433 against 0.6887), with attack recall 0.4975 against 0.3164. The gap
is larger for the models that score best, so the conventional protocol does
not merely inflate a number -- it can recommend the wrong model.
Section~\ref{sec:operational} qualifies what that inversion does and does not
imply.

\resfig{fig_optimism_gap}{The same pipeline, hyperparameters and data
evaluated two ways. Filled markers are the random split, hollow markers the
day-disjoint split; the bar between them is the optimism of the conventional
protocol. The ordering is not preserved -- the weakest detector under the
random split is the strongest once attack families are held out.}{fig:gap}
\resfig{fig_confusion}{Row-normalised confusion for XGBoost, averaged over
seeds, with raw counts in parentheses. The false-positive rate improves under
the day-disjoint protocol while false negatives rise by two orders of
magnitude: the detector does not become noisier, it becomes blind.}
{fig:confusion}

One caveat bounds the comparison: the day boundary removes near-duplicate
flows and attack families at once, and this experiment cannot separate their
contributions. The next section addresses the family component directly.

\subsection{P3: Per-Family Behaviour and Unseen Attack Types}
\label{sec:p3}
Aggregate scores average over families that behave very differently, weighted
by flow count. Figure~\ref{fig:family} breaks the random-split result down by
label, and two things stand out.

First, a macro-F1 of 0.988 coexists with web-attack recalls near zero: on the
same split, XGBoost detects SQL injection at 0.000 and cross-site scripting
at 0.023. These families are a few hundred flows each, so they cost the
aggregate almost nothing -- and they are not the families a defender would
choose to miss. Second, the model that looks worst in aggregate is the only
one that detects them: Logistic Regression reaches 1.00 on XSS and 0.93 on
web brute force where the three stronger models sit between 0.02 and 0.09.
The mechanism is its order-of-magnitude higher false-positive rate; whether
that is a virtue depends on the alert budget, which
Section~\ref{sec:operational} makes quantitative. What it is not is visible
in a macro-averaged number.

\resfig{fig_per_family}{Recall by attack family under the random split. The
model with the weakest aggregate score is the only one that detects the
web-attack families at all.}{fig:family}

\paragraph{Attacks never seen in training.}
Leave-one-family-out removes a family from training entirely, then asks the
detector to flag it. Figure~\ref{fig:loo} shows a division rather than a
gradient. Where a sibling variant remains, generalisation holds: holding out
DoS~Hulk or DoS~slowloris still yields recall near 0.88, other
denial-of-service variants remaining in training. Where the family stands
alone it collapses: botnet traffic is detected at 0.000 and port scanning at
0.0036 across 26{,}822 held-out flows. A supervised flow classifier of this
kind interpolates within a family and does not extrapolate to a new one --
and a detector that misses 99.6\,\% of port scanning when it was absent from
training is not performing anomaly detection in any useful sense, whatever
its benchmark score.

\resfig{fig_leave_one_out}{Leave-one-family-out: recall on a family removed
from training altogether. Detection survives when a related variant remains
and collapses when the family stands alone.}{fig:loo}

\subsection{Feature Selection: Reproducible but Not Canonical}
\label{sec:featstab}
Papers that select the top-$k$ features usually print the list, and a reader
may take it as a statement about which flow properties carry signal. Across
seeds the selection is highly reproducible -- mean pairwise Jaccard 0.756,
with 15 of the 20 features chosen by every seed. Across \emph{criteria} it is
not: impurity and permutation importance, computed on the same split with the
same forest, agree on only 4 of their top 20. Impurity favours packet-length
statistics; permutation puts destination port, forward inter-arrival minimum
and the initial TCP window sizes first. Mutual information agrees closely
with impurity (0.75) and barely with permutation (0.25).

The divergence is the known bias of impurity importance towards
high-cardinality correlated continuous features. It matters here for a
specific reason: no aggregate score would have revealed it, since a pipeline
built on either list reports a similar macro-F1. We therefore present the
selected set as one defensible choice rather than as the informative subset,
release it in full, and note that a paper reporting only an impurity list is
reporting a property of the criterion as much as of the traffic.

% [trimmed for page limit] values are given in the text
%\inputtable{feature_stability}

\subsection{Independent Replication, and P4: Cross-Dataset Transfer}
\label{sec:transfer}
Two experiments are commonly conflated under the heading of cross-dataset
evaluation, as an earlier version of this work conflated them. They answer
different questions.

\emph{Independent-dataset replication} re-runs the whole protocol
\emph{within} UNSW-NB15 -- same pipeline, same hyperparameters, different
corpus (Table~\ref{tab:unsw}). It checks whether a pattern is an artefact of
one benchmark. It is not evidence that a detector fitted on one corpus works
on another, since no model leaves its own dataset.

\emph{Transfer} fits on one corpus and evaluates on the other with no
target-domain data. The datasets share no schema, so we align the seven
flow-level quantities whose semantics match -- duration, forward and backward
packet counts, byte counts and mean packet sizes -- rescaling CICIDS2017
duration from microseconds to seconds. That restriction is itself a confound,
so we measure it: an ordinary in-domain holdout on exactly these seven
features gives XGBoost macro-F1 0.9911 on CICIDS2017 and 0.9289 on UNSW-NB15,
against 0.9879 and 0.9507 on twenty features. The reduced feature space costs
almost nothing in-domain, so what follows is domain shift, not feature
starvation.

Transfer fails in both directions, and differently in each
(Figure~\ref{fig:transfer}). Fitted on CICIDS2017 and tested on UNSW-NB15,
the models call almost everything benign: attack recall is 0.0000 for Random
Forest, 0.0001 for XGBoost, 0.0229 for the MLP. Fitted on UNSW-NB15 and
tested on CICIDS2017 they do the opposite, with false-positive rates from
0.55 to 0.89 and precision between 0.11 and 0.21 despite recall as high as
0.85. Each failure runs towards the majority decision of its source corpus --
CICIDS2017 is 17.7\,\% attack, UNSW-NB15 68.1\,\% -- so what transfers is a
threshold calibrated to one dataset's prevalence and scaling, not a notion of
malicious traffic. Macro-F1 lands between 0.24 and 0.42 against 0.99
in-domain: none of these transferred models is usable.

% [trimmed for page limit] values are given in the text
%\inputtable{unsw}
% [trimmed for page limit] \resfig{fig_transfer}

The conclusion is deliberately narrow. A flow-statistics classifier of this
family, trained on one public corpus, does not carry to another, which
contradicts the practice of citing an in-domain score as evidence of general
applicability. It does not show cross-domain detection to be impossible;
domain adaptation, recalibration on target traffic, and features chosen for
invariance are untested here and are the obvious next step.

\subsection{What the Error Rates Mean in Operation}
\label{sec:operational}
Every number so far was measured on a corpus that is 17.7\,\% attack; a
production link is not. Because precision depends on prevalence while the
error rates do not, the measured rates carry to a realistic prevalence
through Equation~\ref{eq:ppv} with no further experiment.

All four models collapse, and the collapse is governed by the false-positive
rate rather than by the recall the benchmark tables emphasise. Take the most
favourable case in this paper: XGBoost under the random split, recall 0.9935,
false-positive rate 0.0073. At an attack prevalence of one in a thousand its
precision is 0.120 -- seven of every eight alerts false -- and it raises
7{,}257 false alerts per million flows. At one in ten thousand, precision is
0.013.

\inputtable{operational}
\resfig{fig_ppv_prevalence}{Precision against attack prevalence, computed
from the measured error rates. Benchmark corpora sit at the left of this
axis and production links at the right.}{fig:ppv}

\paragraph{A trap for anyone watching alert volume.}
Moving from the random split to the day-disjoint one, the alert burden of
XGBoost \emph{improves}: false alerts per million fall from 7{,}257 to
3{,}006, and precision at one in a thousand barely moves, 0.120 to 0.095. An
operator monitoring alert volume and spot-checking precision would see a
system that had got quieter and no less accurate. What actually happened is
that recall fell from 0.9935 to 0.3164. The degradation is invisible in
exactly the signals a deployment has available -- an argument for holding out
attack families during evaluation rather than hoping to notice later.

\paragraph{The inversion does not recommend Logistic Regression.}
Its false-positive rate is an order of magnitude above the ensembles', so at
one-in-a-thousand prevalence its precision is 0.0083 and it raises 59{,}143
false alerts per million -- some twenty times the burden of XGBoost for a
recall advantage of 0.18. The conclusion is narrower: the conventional
protocol cannot be trusted to rank detectors, and neither can macro-F1 once
prevalence is accounted for. The replication on UNSW-NB15 sharpens this: the
ensembles reach recall above 0.96 there, but at a false-positive rate near
0.075, so precision at one in a thousand is 0.013 rather than 0.120. The
operational value implied by an in-domain score differs by an order of
magnitude between two standard benchmarks.

\subsection{Computational Cost}
\label{sec:cost}
Table~\ref{tab:cost} reports training and amortised inference cost, measured
after a discarded warm-up pass as the mean of five repeats, on a 16-core AMD
workstation with 27.8~GiB of memory running Python 3.13, scikit-learn 1.5.2,
XGBoost 3.4.1 and imbalanced-learn 0.14.2. The two ensembles are
statistically indistinguishable in quality and differ materially in price:
XGBoost trains in roughly a sixth of the time Random Forest needs. Since
Section~\ref{sec:paired} found no accuracy grounds for preferring either,
cost is a legitimate basis for the choice -- and it is the only one this
paper found.

The per-sample figure is model cost amortised over a full batch on all cores.
It is a lower bound on the cost of serving a flow and we do not convert it
into a throughput; what a service sustains is measured in
Section~\ref{sec:serving}.

\inputtable{cost}

\section{Deployment and Measured Serving Cost}
\label{sec:serving}
An offline score says what a detector would decide; it does not say what the
system around it costs to run. This section closes that gap by measuring the
deployed service rather than describing it.

\subsection{The Service and the Load Protocol}
The selected model is exported with the per-column standardisation statistics
of the features it uses and wrapped in a stateless HTTP service implemented
with FastAPI. It exposes a liveness endpoint, a readiness endpoint that
reports failure until the model is resident in memory, a metadata endpoint,
and two scoring endpoints -- one flow and one batch. Every response carries a
header with the in-process scoring time, letting the client separate model
cost from queueing, parsing and transport. The image is pinned, runs as a
non-root user, and runs one worker per replica so that pod-level scaling is
not confounded by in-process parallelism. Figure~\ref{fig:arch} summarises
the architecture. The artefact served is the artefact evaluated: same
pipeline, same training split, same leakage controls, so the latency below
belongs to a model whose detection quality is reported above.

Load is generated in closed loop -- each of $N$ clients issues its next
request as soon as the previous response returns, the behaviour of a bounded
worker pool, and unlike an open-loop rate test it cannot mask saturation
behind an unbounded queue. Each concurrency level begins with a discarded
warm-up phase. Payloads are held-out CICIDS2017 flows replayed in sequence,
so the service parses realistic values. We report the distribution from
median to 99.9th percentile, because a mean conceals exactly the tail that
decides whether an inline sensor is usable.

\subsection{Results}
\label{sec:servingresults}
Table~\ref{tab:serving} and Figure~\ref{fig:serving} report the measurement.
A single replica sustains 309~requests/s at eight concurrent clients, with a
median latency of 24~ms and a 95th percentile of 39~ms. Beyond that point
throughput does not merely stop rising, it falls: at 32 clients the service
delivers 188~requests/s with a 95th percentile of 578~ms, and at 128 clients
128~requests/s with a 99th percentile of 6.2~s. No request failed at any
level. This is ordinary congestion collapse for a single-worker service, and
it is the reason the paper reports a distribution and a saturation point
rather than a peak rate.

\inputtable{serving_latency}
\resfig{fig_serving}{Measured latency distribution and sustained throughput
against concurrent clients. Two panels share one x axis rather than
overloading a single plot with two scales. Throughput peaks at eight clients
and declines thereafter while the tail grows by two orders of
magnitude.}{fig:serving}

\paragraph{What the extrapolation would have claimed.}
The amortised model cost measured in Section~\ref{sec:cost} is
0.001056~ms per flow. Inverted, it implies roughly
$9.5 \times 10^{5}$~flows/s. The service sustains 309~requests/s on the
single-flow path: the extrapolation overstates the measured figure by a
factor of about 3{,}000. Batching recovers much of the gap but not all of it
-- at a batch size of 256 the same replica processes 34{,}662~flows/s, still
27 times below the extrapolation, and it pays a 95th-percentile request
latency of 88~ms for the privilege. The residual is everything the amortised
figure omits: HTTP parsing, JSON deserialisation, per-request Python
overhead, scheduling and queueing. The earlier version of this work quoted
$10^{5}$~predictions/s on exactly this basis; it was wrong by more than two
orders of magnitude, and no amount of care in the offline evaluation would
have revealed that.

\paragraph{Batching is a policy choice, not a free win.}
Figure~\ref{fig:servingbatch} shows the trade explicitly. Moving from
single-flow to a batch of 256 multiplies flow throughput by 114 and
multiplies median request latency by 2.0. An inline sensor that must decide
before forwarding cannot spend the latency; an out-of-band pipeline scoring
exported flow records can, and should. The point is that this is a decision
the deployment makes, with measurable cost on each side, and it is not
visible anywhere in an offline model comparison.

\resfig{fig_serving_batch}{The batching trade-off: flows per second and the
95th-percentile request latency as the batch grows.}{fig:servingbatch}

\paragraph{Cold start.}
Measured from process launch to the first successful readiness response, the
service becomes ready in 7.83~s, of which 1.95~s is loading and warming the
model. The remainder is interpreter and framework start-up. This is the cost
a replica pays before it can serve, and it bounds how quickly capacity can be
added; the scheduling and image-pull components that a cluster would add are
not included.

\subsection{What Is Not Measured}
These numbers characterise a single replica under controlled load on the
machine of Section~\ref{sec:cost}. Autoscaling response, cold start under
scheduling pressure, and request loss during pod replacement are properties
of a cluster rather than of a process, and we do not report them. The
deployment manifests -- resource requests and limits, both probes, and a
horizontal autoscaler with explicit scale-up and scale-down windows -- and
the load profile are released so those measurements can be made. We state
their absence rather than estimating them: substituting an extrapolation for
a measurement here is the error this paper argues against.

\section{Limitations}
\label{sec:limitations}
\paragraph{The benchmark is an instrument, not the world.}
CICIDS2017 was captured in 2017 and carries documented labelling and capture
defects~\cite{engelen2021troubleshooting,lanvin2023errors}, correcting which
changes measured performance. Absolute scores on it should not be read as
estimates of field accuracy, and we do not present them so. The quantities we
emphasise are differences between protocols on identical data with an
identical pipeline, which are far more robust to label noise than levels,
since a defect that inflates one protocol inflates the other. What we do not
establish is the size of the gap on a corrected or modern corpus.

\paragraph{Day-disjoint is not the only leakage-resistant split.}
Grouping by five-tuple or splitting by
timestamp~\cite{moczkodan2026transformers} isolates different sources of
optimism, and a factorial over splitting strategies would separate
near-duplicate flows from unseen families more cleanly than the day boundary
can. We chose it because on this corpus it coincides with the family
boundary, and so answers the operationally important question with one
simply-stated protocol.

\paragraph{The transfer experiment carries a confound.}
Transfer uses seven aligned features rather than twenty, so part of the
degradation could be the reduced feature space rather than the domain shift.
Section~\ref{sec:transfer} reports the in-domain result on the same seven
features so the two can be separated.

\paragraph{No tuning; prevalence modelled, not observed.}
All hyperparameters are fixed a priori, which keeps the protocol comparison
clean but means absolute scores are not the best each family can reach, and
leaves open whether tuning would widen the gap -- we would expect tuning
against a random-split validation set to do so. The operational precision
figures follow exactly from Equation~\ref{eq:ppv} given the measured rates,
but assume those rates transfer to a link with a different traffic mix, which
Section~\ref{sec:transfer} suggests is optimistic; read them as an upper
bound.

\paragraph{The serving measurement is single-node.}
Cluster-level behaviour is not measured, as stated above.

\section{Conclusion}
\label{sec:conclusion}
We asked how much of a benchmark intrusion-detection score belongs to the
detector and how much to the protocol. On CICIDS2017, with everything but the
split held fixed, most of the headline number is protocol: attack recall for
the tree ensembles falls from about 0.99 to about 0.31 once no attack family
is shared between training and test, and the model ranking inverts, so the
conventional protocol does not merely inflate a score -- it can recommend the
wrong model.

Three consequences follow for reporting. A single split is not a result: over
five seeds the difference between the leading models is smaller than the
seed-to-seed variation of the weaker two, and a paired test does not separate
it from zero. An aggregate binary score is not one either, since it averages
over families whose detectability differs by an order of magnitude, and the
families that vanish into the average are not the harmless ones. And
precision measured at benchmark prevalence does not survive the move to a
production link -- where, as Section~\ref{sec:operational} shows, the
degradation we measure is invisible to an operator watching alert volume.

The engineering claim is held to the same standard: we report the measured
latency distribution of the running service rather than deriving it from
model cost, and we name the cluster-level properties we did not measure
instead of estimating them. An earlier version of this work made exactly the
extrapolation we now argue against.

None of this says learned intrusion detection does not work; it says the
published numbers answer an easier question than the operational one. The gap
between protocols estimates that difference, is cheap to compute, and belongs
alongside the conventional figure rather than in place of it. Three
directions follow: re-run the protocol on a corrected~\cite{lanvin2023errors}
and on a contemporary corpus; evaluate one-class and anomaly-based detectors,
designed for unseen families, under the same day-disjoint protocol, where
they may overturn the ordering reported here; and complete the cluster-level
serving measurements, since autoscaling under attack-driven load spikes is
precisely the case an intrusion detector must survive.


\ifanon\else
\subsubsection*{Acknowledgements.} % TODO (tuỳ chọn): lời cảm ơn / nguồn tài trợ.
\fi

% Required by ICTA 2026 (placed immediately before References).
\subsubsection*{AI Declaration.}
Generative AI tools (Anthropic Claude) were used to assist with language
editing, LaTeX preparation, and scaffolding of the experiment code. All
experiments were designed and executed by the authors, every reported number
is an actual measured result, and the authors have reviewed and verified all
content and take full responsibility for this paper.

\bibliographystyle{splncs04}
\bibliography{references}

\end{document}
